\documentclass[11pt]{amsart}
\usepackage[T1]{fontenc}
\usepackage[utf8]{inputenc}
\usepackage{amssymb}
\usepackage{array}
\usepackage[margin=1.05in]{geometry}
\usepackage[hidelinks]{hyperref}
\newcommand{\stamptime}{\the\numexpr\time/60\relax:\ifnum\numexpr\time-60*(\time/60)\relax<10 0\fi\the\numexpr\time-60*(\time/60)\relax}
\AddToHook{shipout/foreground}{\put(0.5\paperwidth,-\paperheight+0.55cm){\makebox[0pt]{\tiny\ttfamily unit/doc/\jobname.tex \ --- \ compiled \today, \stamptime}}}
\newtheorem{theorem}{Theorem}
\newtheorem{proposition}[theorem]{Proposition}
\newtheorem{lemma}[theorem]{Lemma}
\newtheorem{corollary}[theorem]{Corollary}
\theoremstyle{definition}
\newtheorem{definition}[theorem]{Definition}
\newtheorem{statement}[theorem]{Statement}
\theoremstyle{remark}
\newtheorem{remark}[theorem]{Remark}
\newcommand{\Om}{\Omega}
\newcommand{\code}[1]{\texttt{#1}}
\newcommand{\cT}{\mathsf{isTrunc}}
\newcommand{\cC}{\mathsf{isConn}}
\newcommand{\refl}{\mathsf{refl}}
\newcommand{\ap}{\mathsf{ap}}
\newcommand{\Fib}{\mathrm{Fib}}
\newcommand{\fib}{\mathrm{fib}}
\newcommand{\fpr}{\mathsf{fpr}}
\newcommand{\tr}[1]{\|#1\|}
\newcommand{\SCt}{\mathsf{SC3}}
\newcommand{\SCf}{\mathsf{SC4}}
\newcommand{\Tri}{\mathsf{Triples}}
\newcommand{\Del}{\mathsf{Deloopable}}
\newcommand{\Unloop}{\mathsf{Unloop}}
\newcommand{\real}{\mathsf{realise}}
\newcommand{\pteq}{\simeq_*}
\newcommand{\Map}{\mathrm{Map}_*}
\newcommand{\lp}{\mathsf{loop}}
\newcommand{\Z}{\mathbb{Z}}
\newcommand{\cls}[1]{[#1]}
\title{The marked mapping theorem for deloopable targets}
\date{28 October 2026}
\begin{document}
\maketitle
\sloppy

\noindent\textbf{Purpose.} Informal proofs, all formalised (Section~4): (a)
statement M of \code{doc/delooping-embedding.pdf} is equivalent to its looped
form (\code{MarkLoop}); (b) the restricted $2$-torsion lemma, from
deloopability alone (\code{TwoTorsion}); (c) how (b) gives M for every pointed
$2$-connected $4$-type without \code{Presentation4.Pres4} (\code{MarkDirect}),
hence $\pi_4(S^3)=\Z/2$ using David Wärn's delooping embedding (\code{Pi4Direct}; not his intended route through a classification of arbitrary maps between triples, see \code{doc/pi4-s3-embedding.pdf}).

\noindent\textbf{Attribution.} The delooping embedding and the sphere strategy
are David Wärn's, as is the delooping machinery (\code{DeloopMaps}); the
classification by triples is Hart--Milner's (Milner's equivalence). The
cross-term analysis (\code{CrossTerm}) and the bilinear classification
(\code{H4Bilinear}, Theorem B) are those of \code{doc/cross-term.pdf} and
\code{doc/h4.pdf}.

\noindent\textbf{Conventions.} As in the development, $\cC\,n$ means
$(n-2)$-connected, $\cT\,n$ means $(n-2)$-truncated; $\SCf$ is the type of
pointed $2$-connected $4$-types, $\Om:\SCf\to\SCt$ looping, $\real(A,B,k)=\Fib k$
the classification, $\Del\,t=(v:\SCf)\times(\Om v=\real\,t)$
(\code{DeloopEmbedding}). $[X,Y]_*$ is the set truncation of $\Map(X,Y)$
(\code{HomotopyClasses.HC}), $c;d$ composition of classes. $e_m:\Om K(A,m+1)\pteq
K(A,m)$ are the Eilenberg--MacLane identifications (\code{Lev.eA}, \code{eB}).
For $k:K(A,2)\to_*K(B,4)$ put $f_k=e_2;k;e_4^{-1}:\Om K(A,3)\to_*\Om K(B,5)$.

\noindent\textbf{Lemma 1} (\code{SourceS3.LoopEquiv}, $c=2$). If $\cC\,4\,X$ and
$\cT\,6\,Y$, then $\lp:\Map(X,Y)\to\Map(\Om X,\Om Y)$ is an equivalence.

\section{Step 1a: the looped form of M}

Under the $S^3$ axiom let $s\in[S^3,BR]_*$ be the class of
\code{SourceS3.Final}, $e=\code{omega-BR}:\Om BR\pteq R$, and
\[
  r=(\lp s);\cls e\ \in\ [\Om S^3,R]_*,
\]
where $\lp$ on classes is \code{Sphere3Trunc.loopC} (no representative of $s$ is
chosen). For $u=(Y,y,..):\SCf$ write $\Om Y=(y=y)$ and
\[
  \mathsf{M}(u):\ s^*:[BR,Y]_*\to[S^3,Y]_*\ \text{bijective};\qquad
  \mathsf{LM}(u):\ r^*:[R,\Om Y]_*\to[\Om S^3,\Om Y]_*\ \text{bijective}.
\]

\begin{lemma}[looping on classes]\label{lem:loopcls}
If $\cC\,4\,X$ and $\cT\,6\,Y$ then $\lp:[X,Y]_*\to[\Om X,\Om Y]_*$ is a
bijection.
\end{lemma}
\begin{proof}
$\lp$ on classes agrees on representatives with the map of set truncations
induced by $\lp$ on pointed maps (\code{loopC-beta}, \code{STrMap.beta}; both
sides land in a set, so they agree everywhere). The latter is a bijection
because $\lp$ is an equivalence (Lemma 1; \code{classes-equiv}).
\end{proof}

\begin{proposition}[L]\label{prop:L}
For every $u:\SCf$, $\mathsf{M}(u)\iff\mathsf{LM}(u)$.
\end{proposition}
\begin{proof}
Consider the four maps
\[
  s^*:[BR,Y]_*\to[S^3,Y]_*,\quad
  \ell=(\lp s)^*:[\Om BR,\Om Y]_*\to[\Om S^3,\Om Y]_*,\quad
  \epsilon=\cls e^*:[R,\Om Y]_*\to[\Om BR,\Om Y]_*,\quad r^* .
\]
(i) $\lp(s;c)=\ell(\lp c)$ (\code{loopC-comp}). (ii) $r^*c=\ell(\epsilon c)$
(associativity of composition of classes). (iii) $\epsilon$ is a bijection
(precomposition by a pointed equivalence, \code{MapFib.pre-pteq}, on classes
\code{pre-equiv}). (iv) $\lp:[BR,Y]_*\to[\Om BR,\Om Y]_*$ and
$\lp:[S^3,Y]_*\to[\Om S^3,\Om Y]_*$ are bijections (Lemma~\ref{lem:loopcls}:
$\cC\,4\,BR$ is \code{BR-conn}, $\cC\,4\,S^3$ is \code{Sphere3Trunc.Simply.s3-conn},
$\cT\,6\,Y$ is part of $u$). By (i), (iv) and 2-out-of-3, $s^*$ is bijective iff
$\ell$ is; by (ii), (iii) and 2-out-of-3, $r^*$ is bijective iff $\ell$ is.
\end{proof}

\noindent\emph{Formalised}: \code{MarkLoop.LoopCls} (Lemma~\ref{lem:loopcls}),
\code{MarkLoop.LMark[S3,base,surf,ump3]}: \code{r}, \code{At[u].sq}, \code{tri},
\code{e-star}, \code{to-ls}/\code{from-ls}, \code{ls-to-r}/\code{r-to-ls},
\code{to-loop}/\code{from-loop}. Nothing about $Y$ beyond $\cT\,6$ is used, and
no presentation.

\section{Step 1b: the restricted $2$-torsion lemma}

\begin{statement}[T]\label{st:T}
Let $t=(A,B,k):\Tri$ be deloopable, $G:K(\Z,2)\to_*K(A,2)$, and let
$h:\Z\to B$ with $\cls{G;k}=\Phi_B(h)$ (\code{H4Proof.h4}; $\Phi_B(h)=
\cls{\kappa;K(h,4)}$). Then $2\,h(1)=0$. In fact $f_{G;k}$ deloops:
$\Unloop f_{G;k}$ holds.
\end{statement}

\subsection{The proposed reduction to $A=\Z$, checked}
The handoff proposes: deloop $G$ to $\hat G:K(\Z,3)\to_*\tr Y_3$ and pull $Y$
back along it. This is correct. (1) $\Om\tr Y_3\simeq\tr{\Om Y}_2\simeq
\tr{\Fib k}_2\simeq K(A,2)$, the last because $\fpr:\Fib k\to K(A,2)$ has
$2$-connected fibres $k\,x=0$ (torsors under $\Om K(B,4)$), so it is the
$2$-truncation; by construction $\Om Y\to\Om\tr Y_3\simeq K(A,2)$ is $\fpr$
transported along $\Om Y\pteq\Fib k$. (2) $\hat G$ exists by Lemma 1
($K(\Z,3)$ is $\cC\,4$, $\tr Y_3$ is $\cT\,5$). (3) $Y'=K(\Z,3)\times_{\tr Y_3}Y$
is $\cC\,4$ (base $\cC\,4$, fibres those of $Y\to\tr Y_3$, $\cC\,5$) and $\cT\,6$.
(4) $\Om Y'\simeq K(\Z,2)\times_{K(A,2)}\Fib k\simeq\Fib(G;k)$, loops commuting
with pullbacks and (1). So $(\Z,B,G;k)$ is deloopable.

\emph{But the reduction does not reduce the difficulty}: what remains,
\[
  \Fib(\kappa;K(h,4))\ \text{deloopable}\ \Longrightarrow\ 2h(1)=0,
\]
is not Lemma N (\code{Necessity.Nec}), which needs a delooping of the
\emph{cocycle} $\kappa;K(h,4)$, i.e.\ an element of $\Unloop f$ ---~the Crux of
\code{doc/delooping-embedding.pdf} at $(\Z,B,\kappa;K(h,4))$. A delooping of the
\emph{fibre} gives no map $K(\Z,3)\to K(B,5)$. The argument below proves the
Crux at $A=\Z$ directly, and it handles $G$ by naturality, so the reduction (and
its truncation and pullback lemmas) is not needed.

\subsection{Ingredients}
Everything below is stated \emph{pointed in the first variable and pointwise
in the second}: no equality of bipointed maps, and no unit coherence in the
second variable, is ever needed (review point 1).

\noindent\textbf{The cross term, one variable at a time.} For
$f:\Om X\to_*\Om W$ and $q:\Om X$ let $\Delta^f_q:\Om X\to_*\Om W$,
\[
  \Delta^f_q(p)=(fp)^{-1}\cdot f(p\cdot q)\cdot (fq)^{-1},
\]
with any pointing (by ($\lambda$0) below only the underlying function matters).
The bipointed cross term of \code{CrossTerm.Cross} is
$\delta_f=e_2^{-1};\code{Dhat};\tau;\code{Cj}$, and for $x:K(\Z,2)$, $q=e_2^{-1}x$,
the underlying functions satisfy
$\mathrm{ev}_x\circ\delta_f=e_2^{-1};\Delta^f_q;e_4$ pointwise
(\code{Cross.delta-def}, \code{parts-def}, \code{tau-formula}, the definition of
\code{Cj}). Properties, all \emph{pointwise}:
\begin{itemize}
\item[($\Delta$0)] $\Delta^0_q(p)=\refl$.
\item[($\Delta$1)] For $L=\lp P$, $P:X'\to_*X$, and $p',q':\Om X'$:
$\Delta^{L;f}_{q'}(p')=\Delta^f_{Lq'}(Lp')$, from $L(p'\cdot q')=Lp'\cdot Lq'$
(\code{PiOne.Induced.loop-mult}). The same for $L=\hat G=e_2;K(a,2);e_2^{-1}=\lp
K(a,3)$ (\code{Lev.loopK}). No pointing is compared, so the unit case
($p'=\refl$) is not needed. (Were it needed, it could not be obtained ``by $J$
on $q'$'', a loop with fixed endpoints; one would first prove the unit
compatibility of \code{loop-mult} for free paths and then apply $f$.)
\item[($\delta$2)] (Lemma 13, \code{Cross.cell-of}, and the cells argument of
\code{DeloopR.unloop-contr}, with $\Z/2$ replaced by $B$.) For $X=K(\Z,3)$,
$W=K(B,5)$: if $\delta_f=0$ then $f=D_p$ for every $p$, hence $\Unloop f$.
\end{itemize}

\noindent\textbf{The invariant.} For $m:\Om X\to_*\Om W$ and $\gamma\in\Om^3X$
($=\Om^2$ of $\Om X$) put $\lambda_\gamma(m)=(\lp^2m)(\gamma)\in\Om^3W$. Then
\begin{itemize}
\item[($\lambda$0)] (Lemma: pointing irrelevance) if $m.1=m'.1$ pointwise, by
$H$, then $\lambda_\gamma(m)=\lambda_\gamma(m')$. Proof, two levels, each
conjugation identified before it is removed. (i) Pointwise
$(\lp m)(r)=\beta^{-1}\cdot(\lp m')(r)\cdot\beta$ for $r:\Om X$, with
$\beta=m'.2^{-1}\cdot H(\refl)^{-1}\cdot m.2\in\Om^2W$ (\code{loop-ap} on both
sides and naturality of $H$: \code{LoopConj}). (ii) Since $\Om W$ is a loop
space, conjugation by $\beta$ on $\Om(\Om W)$ is pointwise the identity
(Eckmann--Hilton), so $\lp m$ and $\lp m'$ have pointwise equal underlying
functions, by a homotopy $H'$. (iii) Applying (i) one level up,
$(\lp^2m)(\gamma)=\beta'^{-1}\cdot(\lp^2m')(\gamma)\cdot\beta'$ with
$\beta'\in\Om^3W$, and conjugation on $\Om^3W$ is trivial (Eckmann--Hilton). This
is the pattern of \code{Necessity.EPre.At.e-pre} (\code{LoopConj.LoopHomotopy},
then EH); it is to be stated and proved as its own lemma.
\item[($\lambda$1)] $\lambda_\gamma(L;m)=\lambda_{(\lp^2L)\gamma}(m)$
(\code{loop-comp} twice); $\lambda_\gamma(m;e)=(\lp^2e)(\lambda_\gamma m)$;
\item[($\lambda$2)] $\lambda_\gamma(0)=\refl$ (\code{loop-const});
\item[($\lambda$3)] (relation with \code{EDef.E}, review point 2) with
$\gamma_0=(\lp e_1^{-1})(g)\in\Om^2K(\Z,2)$ the generator ($g$ the generator of
the circle) and $\hat\gamma_0=(\lp^2e_2^{-1})(\gamma_0)\in\Om^3K(\Z,3)$, for every
bipointed $c$ and $x$:
\[
  (E\,c)(x)=(\lp^2(\mathrm{ev}_x\circ c))(\gamma_0),\qquad\text{hence}\qquad
  (E\,\delta_f)(x)=(\lp^2e_4)\big(\lambda_{\hat\gamma_0}(\Delta^f_{e_2^{-1}x})\big).
\]
To be stated and proved as its own lemma. The first equation is
\code{ENat.E-loop} and \code{L4.ev-ap} ($\mathrm{ev}_x$ commutes with looping),
with $\lp(\code{phi}\,c;\code{Ex}\,x)$ at $g$ rewritten by \code{loop-comp}. For
the second, compose with $e_4^{-1}$ so that both sides are pointed maps into
$\Om K(B,5)$ with the same underlying function ($\mathrm{ev}_x\circ\delta_f$
above), apply ($\lambda$0) there (this is where the conjugations are identified
and removed), and use ($\lambda$1) for $e_2^{-1}$ and $e_4$. Also
$(\lp^2e_4)(\refl)=\refl$ by the pointing lemma of $\lp$ (\code{loop-const}
style).
\item[($\lambda$4)] (Theorem B) $\Theta c=E\,c;j_B:K(\Z,2)\to_*K(B,2)$ and
\code{BilAt.theta-inj}: $\Theta c=\Theta c'$ implies $c=c'$. The type
$\Map(K(\Z,2),K(B,2))$ is a set (\code{cor5-map}: $K(\Z,2)$ is $1$-connected,
$K(B,2)$ a $2$-type), and $\Om K(B,2)$ is connected, so equality there follows
from equality of the underlying functions (pointing irrelevance,
\code{H4Classes.pi-set}).
\end{itemize}

\noindent\textbf{Two facts about $\fpr:\Fib k\to K(A,2)$}, transported along the
pointed equivalences $\epsilon$, $e_2$ to $\pi$ below.
\begin{itemize}
\item[(F1)] $\Om^2\fpr:\Om^2\Fib k\to\Om^2K(A,2)$ is surjective. By
\code{Fibres.omega-fib} twice, $\Om^2\Fib k\pteq\Fib(\lp^2k)$ over
$\Om^2K(A,2)$; the fibre over $x$ is $(\lp^2k)(x)=\refl$, an identity type of
$\Om^2K(B,4)\simeq K(B,2)$, which is $1$-connected; so the fibre is connected,
in particular merely inhabited.
\item[(F2)] (review point 4: stated for \emph{functions}, unpointed) For a family
of $2$-types $T$ over $K(A,2)$, precomposition with $\fpr$ is an equivalence
$\big((x:K(A,2))\to T\,x\big)\simeq\big((w:\Fib k)\to T(\fpr\,w)\big)$, since the
fibres of $\fpr$ are $2$-connected (\code{ConnectedMaps.Elim}). In particular a
function $\varphi:K(A,2)\to T$ into a $2$-type with $\varphi(\fpr\,w)=t_0$ for
all $w$ satisfies $\varphi(x)=t_0$ for all $x$ (take the family $\varphi(x)=t_0$,
a family of $1$-types).
\end{itemize}

\subsection{The proof}
Let $\epsilon:\Om Y\pteq\Fib k$ witness $\Del\,t$ (\code{PtProp} for $\SCt$),
$\pi=\epsilon;\fpr;e_2^{-1}:\Om Y\to_*\Om K(A,3)$, $f'=f_{G;k}$.

\emph{Step 1 ($\pi$ is a loop map).} By Lemma 1 ($\cC\,4\,Y$; $K(A,3)$ is
$\cT\,5\subseteq\cT\,6$), $\pi=\lp P$ for some $P:Y\to_*K(A,3)$.

\emph{Step 2 ($\pi;f_k$ is null).} $\pi;f_k=\epsilon;\fpr;e_2^{-1};e_2;k;e_4^{-1}
=\epsilon;(\fpr;k);e_4^{-1}$, and $\fpr;k$ is merely $0$
(the analogue of \code{MapFib.SourceP.fseq} for any $k$, to be written:
$(x,w)\mapsto w$ and pointing irrelevance). So merely $N:\pi;f_k=0$ in
$\Map(\Om Y,\Om K(B,5))$.

\emph{Step 3 (transfer along $G$).} $G=K(a,2)$ (\code{H2}, Lemma 16), and
$f'=\hat G;f_k$ with $\hat G=e_2;G;e_2^{-1}=\lp K(a,3)$ (\code{loopK}).

\emph{Step 4 (the goal is a proposition, review point 3).} We prove
$\Theta\,\delta_{f'}=\Theta\,0$, an equality in the set
$\Map(K(\Z,2),K(B,2))$ ($\lambda$4). It is a proposition, so we may eliminate into
it the mere path $N$ of Step 2 and a mere $\gamma'\in\Om^3Y$ with
$(\lp^2\pi)(\gamma')=(\lp^2\hat G)(\hat\gamma_0)$, which exists by (F1).

\emph{Step 5 (the invariant vanishes).} Define
$\varphi:\Om K(A,3)\to\Om^3K(B,5)$, $\varphi(y)=\lambda_{\gamma'}(\pi;\Delta^{f_k}_y)$.
For $b:\Om Y$,
\[
\varphi(\pi b)=\lambda_{\gamma'}(\Delta^{\pi;f_k}_b)
  =\lambda_{\gamma'}(\Delta^{0}_b)=\lambda_{\gamma'}(0)=\refl
\]
by ($\lambda$0) with ($\Delta$1) at $L=\pi=\lp P$, $\ap$ of
$g\mapsto\lambda_{\gamma'}(\Delta^g_b)$ along $N$, ($\lambda$0) with ($\Delta$0),
($\lambda$2). \emph{This is the only use of
deloopability}: the multiplicativity of $\pi$. The target $\Om^3K(B,5)\simeq
K(B,2)$ is a $2$-type, so by (F2) (along $\epsilon$, $e_2$) $\varphi(y)=\refl$
for all $y$. Now for $x:K(\Z,2)$, $q=e_2^{-1}x$:
\begin{align*}
\lambda_{\hat\gamma_0}(\Delta^{f'}_q)
 &=\lambda_{\hat\gamma_0}(\hat G;\Delta^{f_k}_{\hat Gq})
  =\lambda_{(\lp^2\hat G)\hat\gamma_0}(\Delta^{f_k}_{\hat Gq})
  &&\text{($\lambda$0), ($\Delta$1) at $L=\hat G$, ($\lambda$1)}\\
 &=\lambda_{(\lp^2\pi)\gamma'}(\Delta^{f_k}_{\hat Gq})
  =\lambda_{\gamma'}(\pi;\Delta^{f_k}_{\hat Gq})=\varphi(\hat Gq)=\refl
  &&\text{choice of $\gamma'$, ($\lambda$1).}
\end{align*}
By ($\lambda$3), $(E\,\delta_{f'})(x)=(\lp^2e_4)(\refl)=\refl=(E\,0)(x)$
(\code{ETwo.e-zero}), so $\Theta\,\delta_{f'}$ and $\Theta\,0$ have equal
underlying functions, hence are equal ($\lambda$4). By \code{theta-inj},
$\delta_{f'}=0$. (These last two steps are the body of
\code{Necessity.ThetaEq.theta-eq}; it cannot be applied as it stands, since its
hypothesis, the pointwise equality of $E$, is not a proposition and so cannot
receive the mere $N$ and $\gamma'$. We apply its components, \code{pi-set} with
\code{BL.sK2}, \code{BL.cK2}, and \code{theta-inj}, after the elimination; or
split \code{theta-eq} at \code{Theta}.)

\emph{Step 6.} By ($\delta$2), $\Unloop f'$. The goal $2h(1)=0$ is a
proposition and $G;k=\kappa;K(h,4)$ merely (effectiveness of
$\cls{G;k}=\Phi_B(h)$), so we may transport to $\Unloop f_{\kappa;K(h,4)}$, and Corollary conv
(\code{Necessity.Conv}, i.e.\ Lemma N) gives $2h(1)=0$. (The orientations agree:
$f_{\kappa;K(h,4)}$ is \code{Conv}'s \code{CJ.conj F}, conjugation by $e_2$ of
$\Z$ and $e_4$ of $B$; \code{Cross} uses the locked \code{Lev.eB}, equal to
\code{KK.e} by \code{eA-def}/\code{eB-def}.)\qed

\begin{remark}
Taking $G=K(\mathrm{id},2)$ at $A=\Z$, Steps 1--6 prove the Crux of
\code{doc/delooping-embedding.pdf} at every triple $(\Z,B,k)$: if $\Fib k$
deloops, so does $f_k$. For general $A$ the same argument would need Theorem B
with source $K(A,2)$, which the development does not have; Statement T only
needs sources $K(\Z,2)$.
\end{remark}

\begin{remark}[classically]
Steps 3--5 say: the cross term of $k$, a class in $\tilde H^4(K(A,2)\wedge
K(A,2);B)$, vanishes after pulling back to $\Fib k\wedge\Fib k$, because $\Fib k$
is an $H$-space over $K(A,2)$; and pulling back is injective on the
$\pi_2\otimes\pi_2$ part, which is what $E$ detects. The cross term of
$\kappa;K(h,4)$ is $2h\,\iota\otimes\iota$; we do not need this computation,
since Lemma N is used instead.
\end{remark}

\section{Step 1c: M for deloopable targets, and compatibility with $r$}

Let $u=(Y,..):\SCf$ and $t=(A,B,k)$ with $\Om Y\pteq\Fib k$ (\code{sc4-equiv}:
every $u$ arises so). By Proposition~\ref{prop:L} and transport along $\Om
Y\pteq\Fib k$ (\code{ptEqInd}, as \code{rep-pteq}), M at $u$ follows from
\[
  r^*:[R,\Fib k]_*\to[\Om S^3,\Fib k]_*\ \text{bijective}.
\]
Compose with $q=\fpr_*$ on both sides: $q(r;c)=r;q(c)$ (associativity), so the
square commutes. Then:
\begin{enumerate}
\item $r^*:[R,K(A,2)]_*\to[\Om S^3,K(A,2)]_*$ is bijective: Proposition~L at
$u=K(A,3):\SCf$, where M holds (\code{SourceS3.Final.right}), then $e_2$.
\item $q$ is bijective from the source $R$ (\code{MapFib.Lifting}-type lemma for
$k$, with $\Om K(B,4)$ and $\Om\Fib k$ connected): (E) every $c:R\to_*K(A,2)$
is $p_R;G$ on classes ($p_R^*$ bijective into $K(A,2)$ for $R=\Fib\bar k$: a new
lemma, the proof of \code{MapFib.SourceP.PStar} with $\kappa$ replaced by
$\bar k$, the fibres of $p_R$ being $2$-connected), and $\cls{p_R;G;k}=
\cls{p_R;\bar k;K(g,4)}=0$ because $2h(1)=0$ (Statement T), $h=\rho;g$
(\code{Factor}), $\Phi(\rho;g)=\cls{\bar k;K(g,4)}$ (\code{Psi}), and
$p_R;\bar k$ is merely null (the generic \code{fseq} of Step 2); (U) $[R,K(B,3)]_*=0$ (\code{MapBR.UR},
\code{FibreH3.FibAtC} at $(\Z/2,\bar k)$).
\item $q$ is bijective from the source $\Om S^3$: (E) by 1, every class is
$r;c$ with $c$ from $R$, and $\cls{r;c;k}=r;\cls{c;k}=0$ by 2(E); (U)
$[\Om S^3,\Om K(B,4)]_*\cong[S^3,K(B,4)]_*=0$ (Lemma~\ref{lem:loopcls} at $K(B,4)$,
\code{s3-paths}: $\Om^3K(B,4)$ is connected).
\end{enumerate}
By 2-out-of-3, $r^*$ is bijective into $\Fib k$. This proves M for every
$u:\SCf$ with no presentation, and the only new input is Statement T.

\section{Formalisation: status and correspondence}
Everything above is formalised and checked (\code{./segc sag/pi3-s2-postnikov
all}: 96 modules, 0 failed; the 26 \code{.opens} specifications, rule~6 of
\code{seg/WRITING.md}, have no deviation). The correspondence, text $\to$
declaration (modules of \code{sag/pi3-s2-postnikov}):

\smallskip
\noindent
\begin{tabular}{@{}>{\raggedright\arraybackslash}p{0.33\textwidth}>{\raggedright\arraybackslash}p{0.63\textwidth}@{}}
Lemma~\ref{lem:loopcls}, Proposition~\ref{prop:L} & \code{MarkLoop.LoopCls}, \code{MarkLoop.LMark.r}, \code{LMark.At.to-loop}/\code{from-loop}\\
($\lambda$0) & \code{CrossDeloop.PtIrr.lambda0} (\code{LoopIrr}, \code{conj-eh})\\
($\lambda$1), ($\lambda$2) & \code{CrossDeloop.lam-pre}, \code{Inv.lam-zero}\\
($\lambda$3) & \code{CrossDeloop.ELoop.At.e-loop}, \code{EDelta.At.e-delta}\\
($\lambda$4) & \code{H4Bilinear.BilAt.theta-inj}, \code{Necessity.ThetaEq.theta-pt}\\
$\Delta^f_q$, ($\Delta$0), ($\Delta$1) & \code{CrossDeloop.Inv.Delta}, \code{Inv.delta0}, \code{Delta1.d1}\\
$\delta_f$, ($\delta$2) & \code{CrossTerm.CrossGen[B,f].delta}, \code{cell-of}; \code{DeloopR.UnloopGen.unloop-contr}\\
(F1), (F2), generic \code{fseq} & \code{MapFib.FibP.surj2}, \code{FibP.const-ext}, \code{FibP.fseq}\\
Step 1 & \code{TwoTorsion.Loop.eps}, \code{pi}, \code{pi-cell}, \code{pi-loop}, \code{pi-mult}\\
Step 2 & \code{TwoTorsion.Loop.N} (from \code{FibP.fseq}, eliminated in Step~4)\\
Step 3 & \code{TwoTorsion.TwoTors.Ghat}, \code{ghat-loop}, \code{ghat-mult}, \code{f'-eq}\\
Step 4 & \code{TwoTors.tgt}, \code{xA}; the module \code{TwoTors.Elim[n,w,ew]} receives the mere $N$ and $\gamma'$; \code{TwoTors.theta} eliminates them into $\Theta\,\delta=\Theta\,0$ (\code{BL.sK2})\\
Step 5 & \code{Elim.gamma'}, \code{eg}, \code{phi}, \code{phi-pi}, \code{phi-refl}, \code{lam-fp}, \code{h-cell}, \code{e-zero}, \code{theta}\\
Step 6 & \code{TwoTors.cross}, \code{unloop} ($\Unloop f_{G;k}$), \code{TwoTors.Two[h,eh].two-tors} ($2h(1)=0$; \code{conj-bridge}, \code{Necessity.Conv})\\
Section 3, 2(E), 2(U) & \code{MarkDirect.ExR} (with \code{MapFib.FibP[Z,Z/2,$\bar k$].PStar}), \code{ex-R}, \code{un-R} (\code{MapBR.SourceBR.ur}), \code{qR}\\
Section 3, item 1 & \code{MarkDirect.Direct.uA}, \code{markedA} (\code{SourceS3.Final.right}), \code{rA}, \code{transfer}\\
Section 3, item 3 & \code{Direct.ExS}, \code{UnS}, \code{qS}\\
$r^*$ bijective into $\Fib k$; M for every $u$ & \code{Direct.rF}; \code{Direct.AtU[u].marked}, \code{Direct.marked}\\
$\pi_4(S^3)=\Z/2$ & \code{Pi4Direct.Pi4S3.pi4-S3} (\code{Pi4Sphere.Two4}, \code{Pi4Sphere.pi4-BR})
\end{tabular}

\smallskip
\noindent\textbf{Dependencies.} Statement~T (\code{TwoTorsion}) has no sphere
parameter and uses no presentation; its only use of deloopability is
\code{Elim.phi-pi}. \code{MarkDirect} uses Statement~T and the $S^3$ axiom
(through $r$, \code{SourceS3.Final.right} and \code{s3-paths}). The marked class
is the existing one: $r=$ \code{MarkLoop.LMark.r} and $s=$
\code{SourceS3.Final.s}. The new endpoint \code{Pi4Direct.Pi4S3.pi4-S3} does
not depend on \code{Presentation4.Pres4} (nor \code{Final4}) nor on
\code{DeloopEmbedding.Mark.marked}: a copy of the development with
\code{Presentation4.seg} deleted, \code{Mark.At}/\code{Mark.marked} and
\code{Pi4Sphere.Assemble4}/\code{Pi4S3} removed still checks (4874
declarations, 0 failed). \code{Pi4Sphere.Pi4S3} (through \code{Pres4}) is kept
unchanged as the reference.

\noindent\textbf{Changes to existing modules.} \code{MapFib.SourceP.PStar} became
\code{MapFib.FibP.PStar} (Lemma~7 for any triple; \code{SourceP} uses it at
$(\Z,\Z,\kappa)$); \code{Necessity.ThetaEq.theta-eq} was split at $\Theta$
(\code{theta-pt}), as anticipated in Step~5. Two locks were added for rule~6:
\code{Loop.eps} (only its type is used) and \code{Direct.AtU.t} (after its
equation \code{t-eq}); without them a by-definition step opened their
constructions.

\noindent\textbf{Review.} A first review (28 Oct) asked for: ($\Delta$1) only for
actual loop maps, with the pointings explicit; the raw and conjugated
conventions kept apart ($\hat\gamma_0$); the propositional eliminations explicit
(Step 4); (F2) at the level of maps, not classes. The present text follows
these: the argument is pointwise in both variables, the pointings being removed
once, by ($\lambda$0), and it eliminates into the set-level equality
$\Theta\,\delta=\Theta\,0$. A second, independent review checked the revised
text against the code. Its corrections are included: ($\lambda$0) replaces an
incorrect ``by $J$ on $q'$'' for the pointing in ($\Delta$1); the two-level
Eckmann--Hilton step is made explicit; the citations for \code{fseq}, $p_R^*$
and Corollary~C are corrected. A third reading asked that ($\lambda$0) and
($\lambda$3) be separate lemmas with the conjugations identified
(\code{LoopConj}) before Eckmann--Hilton removes them; they were
formalised first (\code{CrossDeloop}).
\end{document}
